\documentclass[11pt,a4paper]{article}

\usepackage[utf8]{inputenc}
\usepackage[T1]{fontenc}
\usepackage{amsmath,amssymb,amsthm}
\usepackage{geometry}
\geometry{margin=1in}
\usepackage{hyperref}
\usepackage{verbatim}
\usepackage{booktabs}
\usepackage{enumitem}
\usepackage{xcolor}

\newtheorem{theorem}{Theorem}[section]
\newtheorem{corollary}[theorem]{Corollary}
\hyphenation{source-intake}

\title{A DASHI Source-Intake Audit of the Eriksson--Ba{\l}aban\\[4pt]
       Yang--Mills Mass-Gap Route}

\author{DASHI Collaboration}
\date{June 2026}

\begin{document}

\maketitle

\begin{abstract}
This manuscript reports a DASHI source-intake audit of the Eriksson--Ba{\l}aban
Yang--Mills route.  The audit records that all mathematical route gates
represented in DASHI close relative to 32 named imported results, conventions,
and source claims: terminal KP certificate, Ba{\l}aban RG consumption, DLR-LSI
discharge, OS0--OS4 Osterwalder--Schrader axioms, OS1/O(4) rotational
covariance, Wightman reconstruction, continuum stability, and OS-route gauge
compatibility.  The result is a conditional, receipt-backed dependency
formalisation, not official Clay recognition.

The promotion authority gate remains closed:
\[
\texttt{clayYangMillsPromoted = false}
\]
because CMI consideration requires publication in a qualifying outlet
(CMI~\S6), a two-year waiting period (\S3(ii)), and broad community
acceptance (\S3(iii)).  None of those conditions is currently satisfied.
\end{abstract}

\section{Introduction and Non-Promotion Statement}

The DASHI repository contains a formally type-checked dependency audit of the
Yang--Mills mass-gap route through the Eriksson--Ba{\l}aban programme
(viXra 2602.0041--2602.0096).  The audit is organised as a network of typed
gate records, each consuming named imported postulates from specific source
papers.

The terminal state of the audit is:

\[
\begin{array}{lcc}
\hline
\text{Field} & \text{Value} & \text{Gate type} \\
\hline
\texttt{mathematicalSourceIntakeClosed} & \texttt{true} & \text{Internal} \\
\texttt{candidateForPeerReview} & \texttt{true} & \text{Readiness} \\
\texttt{candidateForClayTrack} & \texttt{true} & \text{Substantive readiness} \\
\texttt{qualifyingJournalPublication} & \texttt{false} & \text{CMI \S6 --- external} \\
\texttt{twoYearWaitingPeriodElapsed} & \texttt{false} & \text{CMI \S3(ii) --- external} \\
\texttt{globalMathematicsAcceptance} & \texttt{false} & \text{CMI \S3(iii) --- external} \\
\texttt{clayYangMillsPromoted} & \texttt{false} & \text{Authority --- inviolable} \\
\hline
\end{array}
\]

The top three rows are internal to DASHI.  The bottom four rows are external
and cannot be closed by any internal DASHI action.  The CMI's own rules
\cite{clayrules} require a proposed solution to be published in a qualifying
outlet, wait at least two years, and gain general acceptance before
consideration.

The one public claim this repository can make is:

\begin{quote}
\emph{DASHI contains a formally type-checked dependency audit showing that,
conditional on the 32 imported results listed in
\texttt{PostulateInventory.agda}, four-dimensional SU(N) Yang--Mills theory on
$\mathbb{R}^4$ admits a Wightman-axiom-satisfying quantum field theory with a
positive mass gap $\Delta_{\mathrm{phys}} \geq c_N \Lambda_{\mathrm{YM}} > 0$.
This is a conditional source-intake formalisation.
\texttt{clayYangMillsPromoted = false}: official Clay recognition requires
qualifying peer-reviewed publication, a two-year waiting period, and broad
community acceptance, none of which is currently satisfied.}
\end{quote}

\section{Finite Gate 3 Geometry}

The finite geometric foundation is carried by the Yang--Mills Gate 3 receipt
surface.  This section is harvested from the pre-terminal Paper 3 draft and
remains unchanged: the finite carrier geometry is a positive result at every
epoch.

\begin{theorem}[Finite-depth non-flat Yang--Mills curvature witness]
\label{thm:gate3}
The repository contains an inhabited finite Gate 3 receipt,
\texttt{canonicalGate3DiscreteGeometryReceipt} in
\texttt{DASHI.Physics.Closure.YangMillsGate3DiscreteGeometryReceipt}.
This receipt records:

\begin{enumerate}
\item \textbf{Finite discrete forms} (\texttt{canonicalDiscreteForms}):
  $\texttt{siteCarrier} = \texttt{SPTI.ShiftPressurePoint}$,
  $\texttt{zeroFormCarrier} = \texttt{SFGC.SFGCSite2DDiscrete0Form}$,
  $\texttt{oneFormCarrier} = \texttt{SFGC.SFGCSite2DDiscrete1Form}$,
  $\texttt{twoFormCarrier} = \texttt{SFGC.SFGCSite2DDiscrete2Form}$.

\item \textbf{Depth-9 connection surface} (\texttt{canonicalConnectionOnDepth9}):
  $\texttt{connectionWitness} = \texttt{YMObs.sfgcReferenceNonFlat1Form}$,
  $\texttt{curvatureWitness} = \texttt{YMObs.sfgcReferenceNonFlatCurvature2Form}$,
  $\texttt{fieldStrengthWitness} = \texttt{YMObs.sfgcReferenceNonFlatFieldStrengthBridge}$.

\item \textbf{Finite nonzero curvature component}:
  \[
  \texttt{chainB1FiniteFANonzeroAtReference} :
  \texttt{SFGC.sfgcSite2DEvaluate2}
    \texttt{YMObs.sfgcReferenceNonFlatCurvature2Form}
    \texttt{YMObs.sfgcReferencePlaquette}
  \equiv \texttt{SPTI4.phi1}
  \]
  with non-vacuum separation witness
  $\texttt{chainB1FiniteFANotVacuumAtReference}$.

\item \textbf{Local finite nonabelian algebra}:
  $\texttt{bracket} = \texttt{YMObs.localFiniteLie3Bracket}$,
  $\texttt{antisymmetry} = \texttt{YMObs.localFiniteLie3BracketAntisymmetry}$,
  $\texttt{Jacobi} = \texttt{YMObs.localFiniteLie3JacobiWitness}$.

\item \textbf{Local operator law} $D_A^2 = [F_A,\cdot]$ on 0-forms:
  $\texttt{YMObs.localFiniteLie3DASquaredEqualsBracketFAOn0Forms}$.
\end{enumerate}
\end{theorem}

The quantitative content of Theorem~\ref{thm:gate3} is the finite depth-9
datum:

\[
\begin{array}{ll}
\hline
\text{Field} & \text{Value} \\
\hline
\text{Carrier depth} & 9 \\
\text{Two-cell carrier} & \texttt{SFGC.SFGCSite2DPlaquette} \\
\text{Selected plaquette} & \texttt{YMObs.sfgcReferencePlaquette} \\
\text{Evaluated curvature} & \texttt{SPTI4.phi1} \\
\text{Vacuum comparator} & \texttt{SPTI4.phi0} \\
\hline
\end{array}
\]

The depth-9 carrier curvature 2-form is total over all available plaquettes
(\texttt{depth9FiniteCurvatureComponentCoversEveryAvailable2Cell}).  This is a
finite receipt theorem, not a continuum principal-bundle curvature tensor.

\section{Prior Obstacles and Their Resolution}

In the terminal source-intake architecture, the earlier DASHI blockers
\texttt{BalabanPhysicalBetaBridge} and
\texttt{PAWOTGUniformSeparation} are no longer terminal open gates:
their roles are absorbed by the imported Eriksson--Ba{\l}aban source
claims --- specifically DLR-LSI (P12--P15), RG-Cauchy summability
(P16--P18), the all-diameter KP certificate (P06--P11), and the continuum
OS1/Wightman chain (P28--P31).  The resolution is source-intake relative:
DASHI imports these results as named postulates; it does not independently
reprove the underlying analytic estimates.

\subsection{The Beta Gap}

The physical lattice normalisation at $\beta \sim 6$ was KP-divergent:

\[
\begin{array}{lcc}
\hline
\beta & r(\beta) & \text{Status} \\
\hline
6.0   & 2.70 & \text{Divergent} \\
10.11 & 0.99 & \text{Convergence boundary} \\
12.97 & 0.50 & \text{Strict absorption} \\
\hline
\end{array}
\]

The gaps to bridge were $4.11$ (convergence) and $6.97$ (strict absorption).
The perturbative bridge was ruled out by an $\exp(150) \sim 10^{65}$ scale
ratio: with one-loop coefficient $b_0 \sim 0.0465$, a perturbative bridge for
the strict gap requires $\exp(6.97 / 0.0465) \sim \exp(150) \sim 10^{65}$.
This demonstrates why the nonperturbative Ba{\l}aban block-spin construction
is required.

\subsection{Resolution via the Eriksson Series}

The Eriksson series resolves both blockers through three mechanisms:

\begin{description}
\item[DLR-LSI (2602.0052).]
  Lemma~6.3 ($\delta_k < \alpha_{\mathrm{blk}}/4$ via polymer decay +
  $p_0 > 0$, using the Yoshida--Guionnet--Zegarlinski criterion) closes the
  gap deferred by Remark~5.3 of (2602.0088).  Theorem~7.1 gives DLR-LSI
  implying Dobrushin--Shlosman complete analyticity and exponential
  clustering.  Corollary~7.3 gives a positive lattice-level spectral gap
  $\Delta_{\mathrm{phys}} \geq m(\beta, N_c, d) > 0$.

\item[RG-Cauchy summability (2602.0072).]
  Assumption A2 (per-link oscillation bound with $2^{-2k}$ irrelevance) is
  discharged by the all-diameter KP certificate.  Theorem~1.3 closes the
  Efron--Stein influence seminorm (B6).  Corollary~5.1 gives Cauchy
  summability: $\sum_k 2^{-2k} < \infty$, establishing the continuum limit as
  a strict Cauchy sequence.

\item[All-diameter KP certificate (2602.0069, 2602.0056).]
  The Step V spatial KP certificate closes polymer activity control at every
  diameter scale, conditional on 11 named postulates.  The corrected
  $\texttt{absorptionConditionSatisfied}$ gate (replacing the category-error
  $\texttt{cStarGreaterThanOne}$) ensures $C^* \cdot p_0(g_k)$ grows without
  bound via asymptotic freedom.
\end{description}

\section{The 32-Postulate Dependency Audit}

The canonical audit surface is the single typed record at

\[
\texttt{DASHI/Physics/YangMills/PostulateInventory.agda}
\]

The postulates are organised into five topological layers:

\subsection{Layer 0 --- Unconditional}

\[
\begin{array}{cll}
\hline
\text{ID} & \text{Name} & \text{Source} \\
\hline
P01 & \texttt{treePathEdgesExist} & \text{Diestel} \\
P02 & \texttt{graphDistMinimality} & \text{Diestel} \\
P03 & \texttt{treePathBoundedByEdgeCount} & \text{Diestel} \\
P04 & \texttt{kappaStrictlyPositive} & \text{CMP 95 Prop.~1.2} \\
P05 & \texttt{kappaNormalisedToOne} & \text{DASHI A2} \\
P06 & \texttt{ImportedPolymerAnimalCountingBound} & 2602.0041, Lem.~5.6 \\
P07 & \texttt{ImportedKPSummabilityBound} & \text{Derived} \\
P08 & \texttt{pZeroPositive} & \text{CMP 122 eq.~1.89} \\
P09 & \texttt{entropyBeatenByFullDecay} & $\beta \geq \beta_0$ \\
P20 & \texttt{AnisotropicSubspaceClassificationTheorem} & 2602.0087, Thm.~3.6 \\
P28 & \texttt{ImportedRotationalWardIdentity} & 2602.0092, Prop.~3.2 \\
P32 & \texttt{TriangularMixingPreventiveLock} & 2602.0096, Thm.~8.5 \\
\hline
\end{array}
\]

\subsection{Layer 1 --- Conditional on Layer 0}

\[
\begin{array}{cll}
\hline
\text{ID} & \text{Name} & \text{Source} \\
\hline
P10 & \texttt{ImportedLargeFieldActivityBound} & 2602.0069, Thm.~8.5 \\
P11 & \texttt{ImportedAbsorptionCondition} & 2602.0056 \S7 \\
P12 & \texttt{ImportedDLRLSIFromPolymerDecay} & 2602.0052, Lem.~6.3 \\
P13 & \texttt{ImportedCrossScaleBound} & 2602.0052, Lem.~5.7 \\
P14 & \texttt{ImportedDLRLSITheorem} & 2602.0052, Thm.~7.1 \\
P15 & \texttt{ImportedLatticeSpectralGap} & 2602.0052, Cor.~7.3 \\
P21 & \texttt{AnisotropyCoeffQuadraticBound} & 2602.0087, Thm.~5.4 \\
P23 & \texttt{TerminalKPBoundVerified} & 2602.0091, Thm.~1.1+1.2 \\
P24 & \texttt{AssemblyMapComplete} & 2602.0091, Thm.~1.3 \\
P29 & \texttt{ImportedSymanzikBreakingDecomposition} & 2602.0092, Prop.~3.4 \\
\hline
\end{array}
\]

\subsection{Layer 2 --- Conditional on Layer 1}

\[
\begin{array}{cll}
\hline
\text{ID} & \text{Name} & \text{Source} \\
\hline
P16 & \texttt{ImportedAssumptionA2FromKPCertificate} & 2602.0072 A2 \\
P17 & \texttt{ImportedB6InfluenceBound} & 2602.0072, Thm.~1.3 \\
P18 & \texttt{ImportedRGCauchySummability} & 2602.0072, Cor.~5.1 \\
P19 & \texttt{ImportedCouplingControlProof} & 2602.0088, Prop.~4.1 \\
P22 & \texttt{InsertionIntegrabilityBound} & 2602.0087, Thm.~6.6 \\
P25 & \texttt{UniformLSIFixedLattice} & 2602.0089, Thm.~A \\
P26 & \texttt{VolumeUniformMassGapFixedLattice} & 2602.0089, Thm.~B \\
P27 & \texttt{ThermodynamicLimitUnique} & 2602.0089, Thm.~C \\
\hline
\end{array}
\]

\subsection{Layer 3 --- Conditional on Layers 0--2}

\[
\begin{array}{cll}
\hline
\text{ID} & \text{Name} & \text{Source} \\
\hline
P30 & \texttt{ImportedOS1EuclideanCovariance} & 2602.0092, Thm.~4.2+Cor.~4.3 \\
\hline
\end{array}
\]

\subsection{Layer 4 --- Terminal Mathematical Sink}

\[
\begin{array}{cll}
\hline
\text{ID} & \text{Name} & \text{Source} \\
\hline
P31 & \texttt{ImportedWightmanReconstructionWithMassGap} & 2602.0092, Thm.~1.1 \\
\hline
\end{array}
\]

All 32 postulates are \texttt{true}.  The dependency graph is acyclic with a
single sink at P31.  The authority layer (journal publication, waiting period,
global acceptance) has no postulate numbers --- these are not mathematical
objects that DASHI can import.

\section{OS-Route Continuum Construction}

The continuum construction proceeds through the Osterwalder--Schrader axioms,
following the three-tier architecture of (2602.0088, \S8).

\subsection{Tier 1 --- Unconditional (OS0, OS2, OS3)}

\begin{itemize}
\item \textbf{OS0 (temperedness):} Banach--Alaoglu theorem.
\item \textbf{OS2 (reflection positivity):} Osterwalder--Seiler Theorem~2.1;
  follows from the Wilson action and positive-definite Haar measure.
\item \textbf{OS3 (bosonic symmetry):} Automatic for gauge-invariant
  observables.
\end{itemize}

\subsection{Tier 2 --- Hypothesis 5.2 Discharged}

Hypothesis~5.2 (uniform LSI for the pure Wilson measure) is now discharged
via the DLR-LSI Dobrushin bypass (2602.0052, Lem.~6.3 $\rightarrow$ Thm.~7.1).

The remaining Tier 2 gates close:
\begin{itemize}
\item \textbf{OS4 (cluster property):} Exponential clustering from DLR-LSI.
\item \textbf{Physical mass gap:} $\Delta_{\mathrm{phys}} > 0$ (2602.0088
  Cor.~1.2).
\item \textbf{Vacuum uniqueness:} 2602.0088, Prop.~8.5.
\item \textbf{Non-triviality:} 2602.0088, Thm.~8.7.
\item \textbf{Coupling control:} 2602.0088, Prop.~4.1.
\item \textbf{Multiscale decoupling:} 2602.0088, Thm.~6.3.
\end{itemize}

\subsection{Tier 3 --- O(4) / OS1 Restoration}

The O(4) rotational covariance gate is closed by three steps from (2602.0092):

\begin{enumerate}
\item \textbf{Lattice Ward identity} (Prop.~3.2, unconditional):
  \[
  \langle S^\eta_n, L_{\mu\nu} f \rangle =
  -\sum_y \eta^4 \int f(x) \langle (\delta_{\mu\nu} s_W)(y) \cdot
  \prod_j \mathcal{O}(x_j) \rangle\,dx + E^\eta_n(f)
  \]
  with $|E^\eta_n(f)| \leq C_n \eta^2 \|f\|_{W^{1,1}}$.  The rotation
  generator acts on test functions and plaquette positions, not on link
  variables; the Jacobian is identically~1.

\item \textbf{Symanzik decomposition} (Prop.~3.4):
  \[
  (\delta_{\mu\nu} s_W)(y) = g_0(\eta)^{-2} \eta^2
  \bigl[ \lambda_{\mu\nu} \mathcal{O}_{\mathrm{aniso}}(y) +
  Q^{\mathrm{O}(4)}_{\mu\nu}(y) + \mathcal{O}(\eta^2) \bigr]
  \]
  with $\lambda_{\mu\nu} \neq 0$ (Appendix~A, using the 1D anisotropic
  quotient from 2602.0087, Thm.~3.6).

\item \textbf{OS1 Euclidean covariance} (Thm.~4.2 + Cor.~4.3):
  RHS $= \mathcal{O}(\eta^2 \log(\eta^{-1})) \to 0$, so
  $L_{\mu\nu} S_n = 0$ in $\mathcal{S}'(\mathbb{R}^{4n})$.
\end{enumerate}

The triangular mixing preventive lock (2602.0096, Thm.~8.5+Cor.~8.6) confirms
the absence of a $d=4$ anisotropic sink in the $W_4$\text{-scalar} gauge-invariant
sector:
\[
Z_{4\leftarrow 6}\bigl(\mathcal{O}^{W_4,\text{aniso}}_6\bigr) \subset
\mathcal{O}^{W_4}_4 \cap \{\text{O(4)-invariant}\}
\]

\subsection{Wightman Reconstruction (2602.0092, Thm.~1.1)}

With OS0--OS4 and OS1 closed, the Wightman reconstruction theorem gives:
\begin{itemize}
\item A Hilbert space $\mathcal{H}$ with unitary Poincar\'e representation
  $U(\Lambda,a)$.
\item A cyclic vector $\Omega$ (the vacuum).
\item Wightman distributions $W_n$ satisfying all axioms.
\item Positive mass gap $\Delta_{\mathrm{phys}} \geq c_N \Lambda_{\mathrm{YM}} > 0$.
\item Non-trivial scattering: $S^c_4 \not\equiv 0$.
\end{itemize}

\section{Gauge Compatibility for the OS Route}

Per (2602.0096, \S7.6), the OS/Wightman route does NOT require smooth-gauge
fixing.  Gribov singularities have codimension $\geq 2$ in the space of
connections $\mathcal{A}$, and in RCD* theory, subsets of codimension $\geq 2$
have zero capacity.  Therefore they are blind spots for the probabilistic
diffusers.

The Ba{\l}aban RG handles gauge non-perturbatively via small-field/large-field
cutoffs (CMP 122), Faddeev--Popov determinant management at the lattice level
(CMP 99), and block-spin RG preserving reflection positivity (CMP 116).

The gate record therefore has:
\[
\begin{array}{lcl}
\texttt{gaugeFixingCompatibilityClosed} & = & \texttt{true} \quad \text{[OS route]} \\
\texttt{alternateBRSTRouteAvailable} & = & \texttt{false} \quad \text{[out of scope]}
\end{array}
\]

\section{Clay Promotion Authority Gate}

The \texttt{ClayPromotionAuthorityGate} record:

\[
\begin{array}{lcl}
\hline
\text{Field} & \text{Value} & \text{Note} \\
\hline
\texttt{mathematicalSourceIntakeClosed} & \texttt{true} & \text{32 postulates} \\
\texttt{candidateForPeerReview} & \texttt{true} & \text{Auditable} \\
\texttt{candidateForClayTrack} & \texttt{true} & \text{Readiness} \\
\texttt{qualifyingJournalPublication} & \texttt{false} & \text{CMI \S6} \\
\texttt{twoYearWaitingPeriodElapsed} & \texttt{false} & \text{CMI \S3(ii)} \\
\texttt{globalMathematicsAcceptance} & \texttt{false} & \text{CMI \S3(iii)} \\
\texttt{clayOrSABConsiderationAvailable} & \texttt{false} & \\
\texttt{clayYangMillsPromoted} & \texttt{false} & \\
\hline
\end{array}
\]

CMI rules require all three conditions before consideration:
\begin{itemize}
\item \textbf{Qualifying publication (\S6):} The Eriksson series is a viXra
  preprint; DASHI is not a qualifying outlet.  Requires publication in a
  peer-reviewed journal of international standing.
\item \textbf{Two-year waiting period (\S3(ii)):} No such clock has started.
\item \textbf{Global acceptance (\S3(iii)):} CMI and Wikipedia still list
  YM/mass gap as open as of June 2026.
\end{itemize}

\section{Conclusion}

DASHI now contains a type-checked source-intake audit of the
Eriksson--Ba{\l}aban Yang--Mills route.  All mathematical route gates close
relative to 32 named imported postulates.  This is a conditional
formalisation, not an official Clay promotion.  The promotion gate remains
closed:
\[
\texttt{clayYangMillsPromoted = false}
\]
because Clay recognition requires external adjudication: qualifying
peer-reviewed publication, the required waiting period, and broad community
acceptance.  The repository is now source-intake complete and ready for
external mathematical review.

\appendix

\section{Scope of the DASHI Formalisation}

DASHI establishes three things:

\begin{enumerate}
\item \textbf{Type-checking of the dependency structure:} the gate records in
  \texttt{PostulateInventory.agda} are verified by Agda's type checker.
  The dependency DAG is acyclic and has a single mathematical sink (P31).

\item \textbf{Logical correctness of the conditional chain:} if the 32 named
  imported results hold, then the stated conclusions follow by the wiring
  encoded in the gate records.

\item \textbf{Audit transparency:} every imported claim names its source paper,
  theorem number, and conditionality tier.  The audit surface is
  machine-readable and diffable.
\end{enumerate}

DASHI does not establish:

\begin{enumerate}
\item[(A)] \textbf{The analytic estimates inside the Eriksson--Ba{\l}aban
  papers.}  DASHI imports these as postulates; it does not reprove the
  polymer activity bounds, LSI estimates, or Ward identity calculations.

\item[(B)] \textbf{The correctness of the Eriksson series itself.}  That is a
  question for mathematical peer review of those papers.

\item[(C)] \textbf{Clay eligibility.}  See Section~7.
\end{enumerate}

\begin{thebibliography}{9}

\bibitem{clayrules}
Clay Mathematics Institute.
\emph{Rules for the Millennium Prize Problems},
\url{https://www.claymath.org/millennium-problems/rules/}.

\bibitem{eriksson0041}
L.~Eriksson (2026).
\emph{Polymer estimates for lattice Yang--Mills theory}.
viXra:2602.0041.
\url{https://ai.vixra.org/abs/2602.0041}

\bibitem{eriksson0052}
L.~Eriksson (2026).
\emph{DLR-LSI and the spectral gap for lattice Yang--Mills}.
viXra:2602.0052.
\url{https://ai.vixra.org/abs/2602.0052}

\bibitem{eriksson0056}
L.~Eriksson (2026).
\emph{Absorption conditions in the Ba{\l}aban RG}.
viXra:2602.0056.
\url{https://ai.vixra.org/abs/2602.0056}

\bibitem{eriksson0069}
L.~Eriksson (2026).
\emph{Ba{\l}aban--Dimock structural package for SU(N) lattice gauge theory}.
viXra:2602.0069.
\url{https://ai.vixra.org/abs/2602.0069}

\bibitem{eriksson0072}
L.~Eriksson (2026).
\emph{RG-Cauchy summability for blocked Yang--Mills observables}.
viXra:2602.0072.
\url{https://ai.vixra.org/abs/2602.0072}

\bibitem{eriksson0087}
L.~Eriksson (2026).
\emph{O(4) covariance restoration: classification, coefficient bounds,
insertion integrability}.
viXra:2602.0087.
\url{https://ai.vixra.org/abs/2602.0087}

\bibitem{eriksson0088}
L.~Eriksson (2026).
\emph{Continuum stability for lattice Yang--Mills: OS axioms and Wightman
reconstruction}.
viXra:2602.0088.
\url{https://ai.vixra.org/abs/2602.0088}

\bibitem{eriksson0089}
L.~Eriksson (2026).
\emph{Uniform LSI, mass gap, and thermodynamic limit for fixed-lattice
Yang--Mills}.
viXra:2602.0089.
\url{https://ai.vixra.org/abs/2602.0089}

\bibitem{eriksson0091}
L.~Eriksson (2026).
\emph{Terminal KP bound and assembly map}.
viXra:2602.0091.
\url{https://ai.vixra.org/abs/2602.0091}

\bibitem{eriksson0092}
L.~Eriksson (2026).
\emph{Rotational Ward identity, O(4) covariance, and Wightman reconstruction
for Yang--Mills}.
viXra:2602.0092.
\url{https://ai.vixra.org/abs/2602.0092}

\bibitem{eriksson0096}
L.~Eriksson (2026).
\emph{Master map: the Eriksson--Ba{\l}aban Yang--Mills architecture}.
viXra:2602.0096.
\url{https://ai.vixra.org/abs/2602.0096}

\bibitem{balabanCMP95}
T.~Ba{\l}aban (1985).
\emph{Renormalization group approach to lattice gauge field theories I}.
Comm. Math. Phys. 95, 17--40.

\bibitem{balabanCMP99}
T.~Ba{\l}aban (1985).
\emph{Renormalization group approach to lattice gauge field theories II}.
Comm. Math. Phys. 99, 75--102.

\bibitem{balabanCMP116}
T.~Ba{\l}aban (1988).
\emph{Renormalization group approach to lattice gauge field theories III}.
Comm. Math. Phys. 116, 1--22.

\bibitem{balabanCMP119}
T.~Ba{\l}aban (1989).
\emph{Renormalization group approach to lattice gauge field theories IV}.
Comm. Math. Phys. 119, 243--286.

\bibitem{balabanCMP122}
T.~Ba{\l}aban (1989).
\emph{Renormalization group approach to lattice gauge field theories V}.
Comm. Math. Phys. 122, 175--202.

\end{thebibliography}

\end{document}
